\section{Implementations that answer the comparison}
\label{sec:learning}

The first comparison should change the training objective while keeping
the representation fixed.  Use $\epsilon=.025$ and $L=160$, and let
$F_\epsilon$ denote the right-hand side of \eqref{eq:baseline}.
An ordinary PINN approximates one solution after its initial condition
has been fixed.  An operator accepts an initial \emph{function} and
predicts its evolution.  Training a coordinate network on several
different initial fields without an input identifying the field does not
define that operator.

The original FNO learns through Fourier multipliers, pointwise channel
maps and nonlinearities; its demonstrations include Burgers, Darcy and
Navier--Stokes \cite[\S\S3--5]{fno}.  PINO adds physics objectives and
distinguishes operator training from later optimization for one input
\cite[\S3]{pino}.  The construction below is our proposed adaptation to
spatial water-wave Dysthe.  Neither paper establishes its performance
for this equation.

\subsection{One distance-conditioned Fourier backbone}

Let $s\in[0,\Delta]$ be distance within a short propagation slab and
$k_m=2\pi m/L$.  Define the known linear propagator by
\[
 \widehat{S_su}_m=e^{ik_m^2s}\widehat u_m,
\]
which solves $v_s=-iv_{\tau\tau}$.  Use the common ansatz
\begin{equation}
 G_\theta(u,\epsilon,s)=S_su+sC_\theta(u,\epsilon,s).
 \label{eq:conditioned-fno}
\end{equation}
Here $C_\theta$ is a 1D FNO correction, with a declared final Fourier
representation.  It receives the full initial envelope, the distance
$s$, and, if retained as a channel, the fixed value $\epsilon$.
Broadcast scalar conditioning across the profile grid or use a smooth
conditioning map.  Fourier transforms act only in $\tau$; automatic
differentiation acts in $s$.  Propagation distance is not periodic.
Use smooth activations and initially double precision for derivative
and residual checks.

Concretely, lift the input channels to features $z_0$ and use blocks
\[
z_{\ell+1}=\sigma\!\left(W_\ell z_\ell+
\F_\tau^{-1}\!\left[R_\ell(k)\F_\tau z_\ell\right]+b_\ell\right).
\]
Here $W_\ell$ mixes channels pointwise, $R_\ell(k)$ mixes Fourier
channels on a chosen finite set of learned modes, and $\sigma$ is
a smooth pointwise nonlinearity. A final channel map and declared
Fourier projection produce the two components of $C_\theta$.

This ansatz enforces $G_\theta(u,\epsilon,0)=u$ for the supplied
representation and supplies the correct linear phase.  The learned
correction can nevertheless be wrong: the linear control must verify
its behavior.  Setting $\epsilon=0$ retains the cubic NLS term and is
not an isolated linear test.  The same lifting, Fourier layers,
conditioning, projection and precision must be used in the supervised
and physics-informed versions.  Architecture changes belong to later
ablations, once this comparison has passed the numerical controls.

Represent $u$ by its real and imaginary parts.  Do not impose
$\widehat u_{-m}=\overline{\widehat u_m}$ on the complex envelope.
Real-valued feature channels have their own Hermitian symmetry, which
does not imply that restriction on $u$.  Use an endpoint-excluded grid
$\tau_j=\tau_{\min}+jL/N$, $0\leq j<N$, and document whether stored
Fourier coefficients use $e^{ik_m\tau}$ or
$e^{ik_m(\tau-\tau_{\min})}$.  Mixing these conventions introduces a
mode-dependent phase error.  Keep FFT normalization, mode ordering and
the Nyquist treatment consistent with the reference solver.

Fit normalization scales using training data alone.  A common scale for
both complex channels preserves the meaning of phase rotations.
Decode to the physical nondimensional envelope before evaluating
$F_\epsilon$; if normalized variables are used directly, transform
every coefficient and derivative consistently.  Per-input normalization
to unit amplitude discards nonlinear strength unless that amplitude is
also supplied and restored.  Raw absolute position is unnecessary for
this translation-invariant periodic equation.  Test translation and
constant-phase equivariance; a generic two-channel network does not
automatically have the latter symmetry.

\subsection{A full residual requires a whole profile}

For a differentiable reconstructed output $v=G_\theta(u,\epsilon,s)$,
define
\begin{equation}
 r_\theta(u,s)=\partial_s v-F_\epsilon(v).
 \label{eq:operator-residual}
\end{equation}
At each sampled $s$, evaluate a complete periodic slice.  Compute
$\rho=|v|^2$ and $\D\rho$ using the multiplier $|k_m|$, with zero
mean mode, and form all derivative nonlinearities.  Sparse independent
points in $(s,\tau)$ cannot determine this nonlocal term.  Automatic
differentiation through the Fourier evaluation allows training, but
does not remove the need to resolve the profile quadrature.

Declare the continuum reconstruction of grid outputs.  A convenient
choice is an explicit final Fourier projection to $|m|\leq K_{\rm out}$,
with the input represented in the same space.  Report its initial
projection error and refine $K_{\rm out}$.  This final bandwidth is
distinct from the number of learned modes in each FNO layer: skip
paths and pointwise nonlinearities can generate higher frequencies
\cite[\S3 and Appendix E]{pino}\cite[\S\S3--4]{stability}.
For this explicitly truncated output, the cubic right-hand side can
reach $3K_{\rm out}$.  A uniform quadrature with $Q>6K_{\rm out}$
resolves the full residual and its squared spatial $L^2$ norm, up to
floating-point error.  Retain these discarded nonlinear modes in the
loss; projecting the residual back to the output band tests only a
Galerkin system.  This bound does not de-alias arbitrary internal neural
nonlinearities.  Refine the network evaluation grid separately.

For an endpoint-only FNO, a finite difference or an RK defect in
distance measures a discrete update.  It is not
\eqref{eq:operator-residual} without a specified differentiable
reconstruction.  Equation~\eqref{eq:conditioned-fno} supplies that
distance dependence; Fourier differentiation then refers to the
declared profile reconstruction, not an unspecified off-grid network.

\subsection{Matched losses and repeated prediction}

Write $\norm{f}_{2,L}^2=L^{-1}\int_{\TT}|f(\tau)|^2\dd\tau$.
Let $u^{\rm ref}(s;u)$ be the qualified reference starting from $u$,
and choose fixed positive scales $U_*,R_*$ from training information.
For a rollout, set
\[
 \widetilde u_0=u,\qquad
 \widetilde u_j=G_\theta(\widetilde u_{j-1},\epsilon,\Delta),
 \quad 1\leq j\leq J.
\]
A common objective is
\begin{align}
 \mathcal J_{\rm data}
 &=\mathbb E_{u,s}
   \frac{\norm{G_\theta(u,\epsilon,s)-u^{\rm ref}(s;u)}_{2,L}^2}{U_*^2},
 \notag\\
 \mathcal J_{\rm PDE}
 &=\mathbb E_{u,s}\frac{\norm{r_\theta(u,s)}_{2,L}^2}{R_*^2},
 \notag\\
 \mathcal J_{\rm roll}
 &=\mathbb E_u\frac1J\sum_{j=1}^J
   \frac{\norm{\widetilde u_j-u^{\rm ref}(j\Delta;u)}_{2,L}^2}{U_*^2},
 \notag\\
 \mathcal J
 &=\mathcal J_{\rm data}+\lambda_r\mathcal J_{\rm PDE}
   +\lambda_q\mathcal J_{\rm roll}.
 \label{eq:operator-loss}
\end{align}
The expectations denote empirical sampling over training inputs and
distances; $J$ is a short rollout length.  Set $\lambda_r=0$ for the
supervised FNO and $\lambda_r>0$ for the physics-informed version,
with otherwise matched settings.  Declare the weights, optimizer,
budget and selection rule.  Residual distances should include points
independent of labeled snapshots.  Begin by qualifying single-slab
predictions, then introduce the same rollout objective in both models.
A physics-only operator is a separate experiment, not the primary
comparison.

One-step accuracy does not control inputs produced by the model itself.
Brandstetter et al. train on such inputs using a two-step pushforward
with gradients stopped through the first step \cite[\S3.1]{pushforward}.
Test this as a named alternative to differentiating through the short
rollout in \eqref{eq:operator-loss}; do not conflate the two procedures.
At evaluation, use predicted endpoints and never reset to the reference.
Spectral filtering or bounded weights can reduce numerical sensitivity
\cite[\S4]{stability}, but global contraction is not an appropriate
default for modulationally unstable dynamics.  Qualification must
retain the reference sideband growth and phase evolution, rather than
merely produce bounded or visually smooth trajectories.

\subsection{Where a PINN or DeepONet adds information}

Use an ordinary PINN first as a per-input diagnostic: on one short slab,
represent $v_\theta(s,\tau)$ periodically, enforce its initial value,
and minimize the full residual.  Its training cost is incurred for
each input.  A sequential PINN carries its learned endpoint into the
next slab; supplying a reference endpoint instead changes the experiment.
Published failure studies support testing short intervals and explicit
optimization controls \cite[\S5]{failures}.  Causal residual weighting
prioritizes earlier parts of an interval during training
\cite[\S3]{causal}; its weights are optimization diagnostics, not
solution-error certificates.  These findings motivate the proposed
slab strategy but do not guarantee it for Dysthe.

After the shared-backbone comparison passes, an optional
function-conditioned spectral DeepONet can test another representation.
An encoder receives the initial Fourier coefficients; a smooth
distance-conditioned decoder returns complex coefficients
$a_m(s;u)$, giving
\[
 v_\theta(s,\tau;u)=\sum_{|m|\leq K_{\rm out}}
 a_m(s;u)e^{ik_m(\tau-\tau_{\min})}.
\]
Hard initial coefficients and the full residual remain applicable.
This is a proposed spectral variant of the branch/trunk idea
\cite{deeponet}; physics-informed DeepONet has demonstrated learning
across initial functions on other PDEs \cite{pideeponet}.
Neither universality nor those experiments establish optimization
success or generalization to arbitrary water-wave inputs.

Final comparisons use unseen initial realizations, independent profile
grids and distances, raw complex errors, derivative norms, and full
residuals.  Check unweighted residuals even if training uses causal or
adaptive weights.  Keep phase/translation-aligned errors secondary.
Defer test-time weight adaptation: PINO explicitly permits it
\cite[\S3.2]{pino}, but it would mix reusable inference with a new
optimization solve.  If later introduced, report its added cost and
accuracy separately.  The implementation described here is a proposal;
no trained model or long-distance accuracy result is asserted.
